\documentclass[11pt,a4paper]{article}
\usepackage[utf8]{inputenc}
\usepackage[T1]{fontenc}
\usepackage[romanian]{babel}
\usepackage[margin=1cm]{geometry}
\usepackage{array}
\usepackage{tabularx}
\usepackage{enumitem}
\setlist{itemsep=2pt,topsep=2pt,parsep=0pt}
\usepackage{graphicx}
\usepackage{setspace}
\usepackage{xcolor}
\usepackage{hyperref}
\usepackage{amssymb}
\setlength{\parindent}{0pt}
\newcommand{\fillline}[1][6cm]{\underline{\hspace{#1}}}
\newcommand{\fillbox}[1][0.4cm]{\framebox[#1]{\rule{0pt}{#1}}}
\newcommand{\checkboxempty}{$\square$}

\begin{document}

\begin{center}
\includegraphics[width=4cm]{assets/barcode-492006.png}
\end{center}

\vspace{0.3cm}

\begin{minipage}[t]{0.6\textwidth}
Nr. de înregistrare \fillline[4cm] / \fillline[3cm]
\end{minipage}
\hfill
\begin{minipage}[t]{0.3\textwidth}
\fbox{\textbf{Model DITL 001}}
\end{minipage}

\vspace{0.25cm}

\textbf{Către}\\
\textbf{MUNICIPIULUI CLUJ -- NAPOCA}\\
\textbf{DIRECȚIA IMPOZITE ȘI TAXE LOCALE}\\
\textbf{SERVICIUL CONSTATARE, IMPUNERE ȘI CONTROL PERSOANE JURIDICE}\\
e-mail: \underline{registratura@primariaclujnapoca.ro}

\vspace{0.25cm}

\begin{center}
\textbf{CERERE DE COMPENSARE\textsuperscript{1)}}
\end{center}

\vspace{0.3cm}

\hspace{0.5cm}\textit{Subscrisa} \underline{ {{denumire_pj}} }, C.U.I/C.I.F.\textsuperscript{*)} \underline{ {{cui}} }, judeţ \fillline[3cm]\\
loc. \fillline[4cm] cod poştal \fillline[3cm] sector \fillline[1cm], str. \fillline[6cm] nr. \fillline[1.5cm], bloc \fillline[1cm] scara \fillline[1cm]\\
etaj \fillline[1cm] ap \fillline[1cm] tel. \fillline[3cm] fax \fillline[3cm], adresă de e-mail \fillline[5cm], înregistrat la\\
registrul comerţului \fillline[4cm] la nr. \fillline[3cm], reprezentată\\
prin \underline{ {{nume_complet}} } în calitate de proprietar/coproprietar/ mandatar/ acţionar unic\\
/asociat/administrator/împuternicit cu domiciliul în \underline{ {{judet}} } \underline{ {{localitate}} },\\
strada \underline{ {{strada}} }, nr. \underline{ {{numar}} }, bl. \underline{ {{bloc}} }, sc. \underline{ {{scara}} }, et. \underline{ {{etaj}} }, ap. \underline{ {{apartament}} }, legitimat(ă) prin B.I./C.I./C.I.P./Paşaport\\
seria \underline{ {{ci_seria}} } nr. \underline{ {{ci_numar}} }, CNP \underline{ {{cnp}} }, în conformitate cu art. 167 alin (7) din Legea nr. 207/2015 privind\\
Codul de procedură fiscală, cu modificările şi completările ulterioare, solicit compensarea sumei de \underline{ {{suma_compensare}} } lei,\\
reprezentând: \underline{ {{tip_creanta}} } cu \fillline[6cm]\\
achitată prin chitanţă/mandat poştal/ordin de plată \hrulefill\\

\hspace{0.5cm}Anexez prezentei următoarele documente \textsuperscript{2} \hrulefill\\
\hrulefill\\
\hrulefill

\vspace{0.3cm}

Prin semnarea prezentei am luat la cunoştinţă ca declararea necorespunzătoare a adevărului se pedepseşte conform legii penale, cele declarate fiind corecte şi complete.

\vspace{0.2cm}

\checkboxempty\ Sunt de acord ca actele administrative fiscale să-mi fie comunicate la adresa de poștă electronică (adresa e-mail)\\
\checkboxempty\ Mă oblig să comunic instituției orice modificare intervine în legătură cu această adresă de e-mail\\
\checkboxempty\ Îmi exprim consimțământul ca Primăria Municipiului Cluj-Napoca să comunice orice informații, date personale, clarificări și completări pe adresa de e-mail indicată mai sus\\
\checkboxempty\ Am luat la cunoștință faptul că în cazul nefuncționării serverului de e-mail comunicat sau în cazul adresei greșite de e-mail, Municipiul Cluj-Napoca nu poate fi tras la răspundere pentru acest lucru

\vspace{0.25cm}

Data \fillline[4cm] \hfill Semnătura \fillline[5cm]

\vspace{0.3cm}
\hrule
\vspace{0.2cm}

\scriptsize
\textsuperscript{1)} Compensarea se constată de către organul fiscal competent, la cererea debitorului sau din oficiu. Dispozițiile art. 165 privind ordinea stingerii datoriilor sunt aplicabile în mod corespunzător.\\
\textsuperscript{*)} se va completa: codul de identificare fiscală (codul numeric personal, numărul de identificare fiscală, codul de înregistrare fiscală sau codul unic de înregistrare, după caz); actul de identitate seria, numărul; numărul de înregistrare la instanţa judecătorească, etc

\vspace{0.3cm}

Timp estimativ de completare: 15 minute

\vspace{0.2cm}

Datele personale care vă sunt solicitate prin prezenta cerere vor fi prelucrate numai în vederea procesării și soluționării solicitării dumneavoastră. Primăria municipiului Cluj-Napoca garantează securitatea procesării datelor și arhivarea acestora în conformitate cu prevederile legale în vigoare. Responsabilul Primăriei municipiului Cluj-Napoca cu protecția datelor poate fi contactat pe adresa de email dpo@primariaclujnapoca.ro.

În conformitate cu Regulamentul nr. 679/2016 aveți dreptul de a solicita Primăriei municipiului Cluj-Napoca, în ceea ce priveşte datele cu caracter personal referitoare la persoana vizată, accesul la acestea, rectificarea sau ştergerea acestora sau restricţionarea prelucrării sau a dreptului de a vă opune prelucrării, precum şi a dreptului la portabilitatea datelor.

\textsuperscript{2)} Primăria municipiului Cluj-Napoca asigură gratuit fotocopierea documentelor. Pentru documentele emise de către alte instituții publice, organe de specialitate ale administrației publice centrale și locale, precum și persoane juridice de drept privat, care potrivit legii au obținut statut de utilitate publică sau sunt autorizate să presteze un serviciu public, în regim de putere publică, pe care solicitantul nu le depune, acesta are posibilitatea de a-și da consimțământul expres ca Primăria municipiului Cluj-Napoca să obțină în numele său copii sau extrase ale acestora.

\end{document}
